If a nervous system's wiring diagram is enough to predict behaviour, that prediction should
not depend on which individual's wiring diagram is used. The whole-brain leaky
integrate-and-fire (LIF) model of Shiu et al.\ \citep{shiu2024} is the strongest test of the
first half of that statement. It was built directly from the FlyWire connectome of one adult
female \textit{Drosophila} \citep{dorkenwald2024,schlegel2024}. Every neuron has the same
parameters, every synapse has the same weight $\Wsyn$, and each neuron's sign comes from its
predicted transmitter \citep{eckstein2024}. With one tuned parameter, the model nevertheless
predicted which neurons drive proboscis extension and antennal grooming, and 91\% of 164
experimentally testable predictions were consistent with experiments \citep{shiu2024}. Such
connectome-constrained models are now used to generate hypotheses at scale
\citep{pospisil2024,lappalainen2024}. Yet the wiring alone does not in general determine
circuit dynamics \citep{bargmann2013,beiran2025}, and every prediction from the Shiu model
rests on a sample of one animal.

Two further whole-CNS connectomes now make the second half of the statement testable. The
brain-and-nerve-cord connectome BANC \citep{bates2026} is a second female. The male CNS
\citep{berg2026} is a male, reconstructed by a different team and pipeline. Most of their
cell types can be matched to FlyWire \citep{schlegel2024,berg2026}. Connectomes of different
individuals share a stereotyped core but differ in detail, both in worms \citep{witvliet2021}
and in flies \citep{schlegel2024,pedigo2023}. Crucially, two connectomes can differ for three
reasons that look identical in the data: true individual variation, differences in
reconstruction (imaging, segmentation, synapse detection, proofreading completeness), and sex.

Independent groups have already started porting the Shiu model. An open project ran it
unchanged on FlyWire, the male CNS (v0.9) and BANC, and found that sugar-receptor stimulation
drives the proboscis motor neuron MN9 at 76.7, 68.9 and 17.8\,Hz respectively
\citep{brainlab2026}. A pre-registered benchmark recovered sugar$\to$MN9 in the male CNS v1.0
after recalibration \citep{fugsukjit2026}. Most directly, Silva Santos \citep{santos2026} (code archive with a draft preprint, not peer reviewed) reimplemented
the model exactly and compared FlyWire with BANC across six stimulation protocols. Cell-type
response profiles agreed strongly between repeated simulations (Spearman $\rho=0.98$),
moderately between the two hemispheres of one fly ($\rho=0.54$), and weakly between the two
flies ($\rho\le0.32$). BANC as released lost sugar$\to$MN9 and JO$\to$aBN1. That study also
showed that BANC's edge lists carry about 2.1-fold fewer synapses per neuron on matched types,
and attributed the gap to synapse detection rather than biology. In a single, much smaller
circuit, differences between reconstructions of the same animal rival differences between
animals \citep{yad2026}.

Whether the FlyWire--BANC disagreement reflects individuality or reconstruction has not been
tested. Nor has a third connectome been placed on the same agreement scale, or the observed
agreement been compared with a null. Here we address these gaps. We (i) validate a batched
GPU implementation of the Shiu model against the original Brian2 code; (ii) run Shiu's
experimentally validated feeding and grooming predictions on FlyWire v783, BANC v888 and the
male CNS v1.0 under three pre-specified weight normalisations; (iii) place all three
connectomes, both hemispheres and two kinds of control on one agreement ladder; and (iv) ask
directly whether removing synapses from FlyWire until each cell type matches BANC's input
deficit (a ``degraded FlyWire'') reproduces BANC's behaviour. Degree-preserving shuffles of
each connectome \citep{maslov2002} set the floor of the ladder. A robustness sweep over
synapse thresholds, normalisations and sign rules, and a stimulation screen over all shared
sensory types, test how general the conclusions are.
